\documentclass[12pt,a4paper,oneside]{report}

\usepackage[utf8]{inputenc}
\usepackage[T1]{fontenc}
\usepackage[spanish,es-nodecimaldot]{babel}
\usepackage{lmodern}
\usepackage{microtype}
\usepackage{geometry}
\usepackage{setspace}
\usepackage{enumitem}
\usepackage{booktabs}
\usepackage{tabularx}
\usepackage{amsmath,amssymb,amsfonts}
\usepackage{xcolor}
\usepackage{tikz}
\usepackage{graphicx}
\usepackage{listings}
\usepackage[pdftex,hidelinks,pdfauthor={Francisco Clarke},pdftitle={Orquestacion de agentes de lenguaje para el desarrollo de software con descomposicion adaptativa}]{hyperref}

\usetikzlibrary{arrows.meta,positioning,shapes.geometric,calc}

\geometry{top=3.0cm,bottom=2.5cm,left=3.0cm,right=2.5cm}
\onehalfspacing

\lstset{
  basicstyle=\ttfamily\small,
  keywordstyle=\color{blue!80!black}\bfseries,
  stringstyle=\color{green!50!black},
  commentstyle=\color{gray}\itshape,
  numbers=left,
  numberstyle=\tiny\color{gray},
  numbersep=8pt,
  showstringspaces=false,
  breaklines=true,
  frame=single,
  backgroundcolor=\color{black!2}
}

\newcommand{\autor}{Francisco Clarke}
\newcommand{\lu}{120547}
\newcommand{\carrera}{Ingeniería en Sistemas de Información}
\newcommand{\departamento}{Departamento de Ciencias e Ingeniería de la Computación}
\newcommand{\universidad}{Universidad Nacional del Sur}
\newcommand{\directora}{Dra. Ana Gabriela Maguitman}
\newcommand{\codirector}{Dr. Federico Martín Schmidt}
\newcommand{\titulo}{Orquestación de agentes de lenguaje para el desarrollo de software: una política de descomposición adaptativa con verificación basada en evidencia}
\newcommand{\sistema}{\textbf{ManyHands}}
\newcommand{\ctask}{C_{\mathit{task}}}

\begin{document}
\hypersetup{pageanchor=false}

% ==============================================================================
% PORTADA
% ==============================================================================
\begin{titlepage}
  \begin{center}
    {\large\textbf{\universidad}}\\[0.4em]
    {\large\textbf{\departamento}}\\[2.5em]

    {\Large\textbf{TRABAJO FINAL DE \carrera}}\\[3em]

    {\Large\textbf{\titulo}}\\[4em]

    \begin{tabular}{rl}
      \textbf{Alumno:} & \autor \ (LU: \lu) \\[0.8em]
      \textbf{Directora:} & \directora \\[0.8em]
      \textbf{Codirector:} & \codirector \\
    \end{tabular}

    \vfill
    {\large Bahía Blanca, Argentina}\\[0.5em]
    {\large 2026}
  \end{center}
\end{titlepage}

\hypersetup{pageanchor=true}
\pagenumbering{roman}

% ==============================================================================
% RESUMEN
% ==============================================================================
\chapter*{Resumen}
\addcontentsline{toc}{chapter}{Resumen}

Los agentes de programación basados en modelos de lenguaje pueden modificar un
repositorio, ejecutar pruebas y reaccionar a sus resultados. Cuando el objetivo
atraviesa varias capas, sin embargo, un único agente concentra demasiado
contexto y una división excesiva introduce coordinación, contratos y conflictos
de integración. Esta tesis estudia esa tensión, denominada \emph{paradoja de la
granularidad}, mediante el diseño, la implementación y la evaluación de
\sistema, un orquestador local de agentes de código.

El aporte técnico combina dos responsabilidades. Un planificador semántico
propone unidades de trabajo justificadas por el objetivo y el repositorio. Una
política determinista y versionada evalúa si conviene conservar cada unidad como
hoja cohesiva o aceptar un corte semántico propuesto. El resultado se compila en
un grafo híbrido con jerarquía de integración y relaciones tipadas para
artefactos, interfaces y conflictos. Cada intento se ejecuta en un
\emph{worktree} aislado; el orquestador inspecciona el diff, controla el alcance,
valida el commit candidato exacto, integra de abajo hacia arriba y publica sólo
un resultado acompañado por manifiesto y comprobante. La historia canónica del
run es un journal append-only de eventos de dominio.

La evaluación confirmatoria final se pre-registró como una serie pequeña y
acotada. Se utilizó un repositorio Node ESM con capas de dominio, aplicación y
API; dos formas de tarea ---una modificación multi-capa y una función cohesiva
de dominio---; dos repeticiones por forma; una única configuración adaptativa;
el modelo Codex \texttt{gpt-5.4-mini}; cero reintentos automáticos; y un oráculo
externo ejecutado sobre cada SHA candidato exacto. Un rehearsal previo verificó
el instrumento y quedó fuera del denominador.

Se fijaron dos hipótesis. \textbf{H-F1} exigía que las cuatro celdas llegaran a
\texttt{completed} y \texttt{delivered}, publicaran un SHA no vacío y pasaran el
oráculo externo. \textbf{H-F2} exigía que la tarea multi-capa produjera una raíz
compuesta con tres hojas ---dominio, aplicación y API---, y que la tarea cohesiva
produjera una sola hoja ejecutable. El resultado fue \textbf{PASS 4/4 para
H-F1} y \textbf{PASS 4/4 para H-F2}. Las cuatro entregas y topologías se derivan
de journals, snapshots, recibos y resultados de oráculo preservados junto con
el manuscrito.

El alcance de la conclusión es deliberadamente estrecho. La serie no tiene
controles A/B y usa un único target pequeño; por ello no demuestra superioridad
frente a no descomponer, causalidad, optimalidad, escalabilidad ni
generalización estadística. Las líneas históricas Warehouse, G5, G6, G7 y SP2
no se combinan con el resultado final: Warehouse permanece en 1/8, G5 produjo un
resultado comparativo negativo, G6 fue inconcluso, G7 no produjo candidatos y
SP2 fue piloto. Se conservan como evidencia adversa o formativa, no como piezas
incompletas de un denominador positivo.

La conclusión defendible es concreta: bajo la versión, el modelo, la política,
los límites y el protocolo congelados, \sistema\ realizó de forma reproducible
el recorrido de planificación, ejecución aislada, integración, validación
externa del commit exacto y entrega, y seleccionó las dos formas de topología
pre-registradas. El trabajo aporta además una disciplina de evidencia para no
confundir actividad de agentes, tests autocontenidos o resultados parciales con
software entregado y verificado.

\textbf{Palabras clave:} orquestación de agentes, modelos de lenguaje,
descomposición de tareas, granularidad adaptativa, worktrees de Git, validación
de commits, evidencia experimental.


\tableofcontents
\listoffigures
\listoftables

\clearpage
\pagenumbering{arabic}

% ==============================================================================
\chapter{Introducción}

\section{Motivación}

La programación asistida por modelos de lenguaje atravesó en pocos años tres
etapas: autocompletado dentro del editor, generación de fragmentos a partir de
una descripción y, finalmente, agentes autónomos capaces de leer un repositorio,
ejecutar comandos y reaccionar a los errores del compilador o de las pruebas
\cite{chen2021codex,yao2023react}. Los resultados sobre incidencias reales de
proyectos abiertos muestran que estos agentes ya resuelven una fracción no
trivial de tareas de mantenimiento \cite{jimenez2024swebench,yang2024sweagent}.

El salto pendiente no es de capacidad del modelo sino de \emph{organización del
trabajo}. Un objetivo como ``agregar categorías de gasto al sistema, exponerlas
en la API y mostrarlas en la interfaz'' no es un fragmento de código: es un
conjunto de cambios coordinados sobre archivos distintos, con contratos
implícitos entre ellos y con la obligación de no romper lo que ya funciona.
Delegarlo a un agente único expone tres patologías bien documentadas en la
práctica:

\begin{enumerate}
  \item \textbf{Saturación de contexto.} El agente intenta sostener demasiados
  archivos simultáneamente; a medida que crece la masa de tokens, aumenta la
  probabilidad de que invente firmas, pierda importaciones o edite código ajeno
  al pedido.
  \item \textbf{Desborde de alcance.} Sin límites explícitos, el agente modifica
  configuración global o componentes compartidos y rompe módulos que no estaban
  bajo revisión.
  \item \textbf{Ausencia de aislamiento.} Ejecutar un agente directamente sobre
  el directorio de trabajo del desarrollador pone en riesgo trabajo no
  guardado.
\end{enumerate}

La respuesta natural es dividir. Pero dividir bien no es evidente, y ese es el
punto de partida de este trabajo.

\section{Planteamiento del problema: la paradoja de la granularidad}

Descomponer un objetivo de software en unidades ejecutables enfrenta una tensión
que la ingeniería de software conoce desde hace décadas. Parnas
\cite{parnas1972criteria} estableció que los módulos deben delimitarse por las
decisiones de diseño que ocultan, no por los pasos del procesamiento; una
división que ignora ese criterio produce módulos que no pueden evolucionar de
forma independiente. Brooks \cite{brooks1995mythical} mostró que agregar
participantes a una tarea incrementa el costo de comunicación de forma
superlineal, de modo que subdividir tiene un costo propio. Conway
\cite{conway1968committees} observó que la estructura del sistema producido
refleja la estructura de comunicación de quienes lo construyen.

Trasladadas a la orquestación de agentes, estas observaciones delimitan dos
fallas simétricas:

\begin{itemize}
  \item \textbf{Sub-división:} una unidad declarada como hoja sigue siendo
  demasiado amplia. El agente que la recibe reproduce exactamente las patologías
  del agente único.
  \item \textbf{Sobre-división:} un objetivo simple se parte en unidades
  triviales. Cada unidad requiere su propio contexto, su propio contrato de
  interfaz y su propia validación; el costo de coordinación y el riesgo de
  conflicto durante la integración crecen sin beneficio.
\end{itemize}

El problema es que el punto intermedio no es fijo. Depende del objetivo, de la
forma del repositorio, de cuántas interfaces públicas se tocan y de cuánta
superficie de validación involucra el cambio. Una regla estática ---``dividir
siempre en cuatro'', ``una tarea por archivo''--- falla en un extremo o en el
otro. \textbf{El problema que aborda esta tesis es cómo decidir esa granularidad
de forma adaptativa, explicable y reproducible.}

\section{Objetivos}

\subsection{Objetivo general}

Diseñar, implementar y evaluar una plataforma local de orquestación de agentes
de lenguaje que adapte la granularidad de un objetivo de software a su forma
semántica y entregue únicamente resultados verificables sobre un commit exacto.

\subsection{Objetivos específicos}

\begin{enumerate}
  \item Separar el juicio semántico que propone un corte de la política
  determinista que decide si aceptarlo y del compilador que lo transforma en un
  grafo ejecutable.
  \item Representar jerarquía, disponibilidad de artefactos, compatibilidad de
  interfaces y riesgo de conflicto mediante relaciones distintas y tipadas.
  \item Aislar cada intento en un worktree, controlar su alcance y adoptar el
  resultado sólo contra sus entradas exactas.
  \item Vincular cada criterio de aceptación con evidencia obtenida sobre el
  candidate SHA, incluyendo baseline, controles negativos e integridad de
  tests.
  \item Integrar de abajo hacia arriba y publicar sólo después de validar el
  candidato final y emitir un recibo de entrega.
  \item Persistir la historia del run como eventos append-only, separando hechos
  de dominio, snapshots y trazas diagnósticas.
  \item Evaluar dos proposiciones acotadas mediante una serie pre-registrada,
  repetida y verificada por un oráculo externo.
\end{enumerate}

\section{Preguntas de investigación e hipótesis}

\begin{description}[leftmargin=2.2cm,style=nextline]
  \item[\textbf{PI-1}] ¿Puede ManyHands llevar cambios de distinta forma desde
  el objetivo hasta una entrega confirmada cuyo commit exacto pase un oráculo
  externo?
  \item[\textbf{PI-2}] ¿Selecciona la política una topología multi-hoja para una
  tarea que atraviesa fronteras de integración y una única hoja para una tarea
  cohesiva?
\end{description}

PI-1 se operacionaliza como H-F1: cuatro celdas \texttt{completed} y
\texttt{delivered}, cuatro SHAs no vacíos y cuatro oráculos PASS. PI-2 se
operacionaliza como H-F2: dos topologías de tres hojas para la tarea multi-capa
y dos topologías de una hoja para la tarea cohesiva. Las reglas completas se
presentan en el capítulo~\ref{cap:evaluacion}.

Estas preguntas no incluyen superioridad frente a no descomponer. Tampoco
preguntan si la topología observada es globalmente óptima. El diseño final elige
proposiciones que el instrumento puede medir de forma completa, en vez de
mantener hipótesis históricas que quedaron fuera del régimen o no produjeron
candidatos.

\section{Alcance y limitaciones}

\sistema\ es un plano de control local y de un solo usuario. Planning,
scheduling, persistencia, integración y validación se ejecutan en la máquina del
desarrollador; la inferencia se obtiene mediante CLI de proveedores remotos. No
se afirma privacidad absoluta del código enviado al modelo.

La evaluación confirmatoria usa un target Node ESM pequeño, un modelo, una
política y dos repeticiones por forma de tarea. Permite estudiar viabilidad y
trazabilidad end-to-end, pero no escalabilidad, superioridad, causalidad,
optimalidad ni generalización estadística.

Las series históricas no se borran: Warehouse queda 1/8, G5 negativo, G6
inconcluso, G7 sin candidatos y SP2 como piloto. Se reportan como antecedentes
adversos o formativos, pero no se combinan con las cuatro celdas finales.

Quedan fuera del alcance: selección adaptativa de modelos, soporte multiusuario,
aislamiento por contenedores, calibración empírica de pesos, generalización a
otros lenguajes y auditoría externa de accesibilidad. También se declara como
gap de transición la materialización selectiva de artefactos por archivo: el
contrato puede expresarla, pero el driver actual aún transporta el commit
producido completo.


\section{Organización del documento}

El capítulo~\ref{cap:preliminares} introduce los conceptos necesarios. El
capítulo~\ref{cap:estado} revisa el estado del arte y posiciona el trabajo. El
capítulo~\ref{cap:metodologia} presenta la política de descomposición
adaptativa, que constituye el aporte central. Los capítulos~\ref{cap:arquitectura}
y~\ref{cap:implementacion} describen la arquitectura y su implementación. El
capítulo~\ref{cap:evaluacion} reporta la evaluación. El
capítulo~\ref{cap:discusion} discute los resultados, las amenazas a la validez y
las limitaciones. El capítulo~\ref{cap:conclusiones} concluye.

% ==============================================================================
\chapter{Conceptos preliminares}
\label{cap:preliminares}

\section{Agentes de código basados en modelos de lenguaje}

Un \emph{agente de código} es un sistema que envuelve a un modelo de lenguaje en
un bucle de razonamiento y acción: el modelo propone una acción (leer un
archivo, ejecutar una prueba, aplicar una edición), un ejecutor la realiza, y el
resultado se devuelve al modelo como observación \cite{yao2023react}. Esta
arquitectura permite al modelo reaccionar a errores de compilación o a pruebas
fallidas sin intervención humana \cite{shinn2023reflexion}.

En este trabajo, los agentes son ejecutores de línea de comandos de nivel
profesional (Codex CLI, Claude Code) que ya implementan ese bucle. \sistema\ no
los reemplaza: los \emph{orquesta}, proveyéndoles el aislamiento, el alcance
acotado y los contratos que necesitan para operar sobre un repositorio real.

\section{Grafos dirigidos acíclicos de tareas}

Representar un plan de ejecución como un grafo dirigido acíclico (DAG) es una
técnica consolidada. En el contexto de modelos de lenguaje, LLMCompiler
\cite{kim2024llmcompiler} organiza llamadas a herramientas en un DAG cuyas
aristas son dependencias de datos, lo que permite identificar y ejecutar en
paralelo las ramas independientes.

Aplicado al desarrollo de software, el modelo requiere aristas más ricas que
``el resultado de A alimenta a B''. Este trabajo distingue cuatro relaciones
(sección~\ref{sec:relaciones}) porque la jerarquía de integración, la
disponibilidad de un artefacto, la compatibilidad de una interfaz y la exclusión
por recurso compartido tienen consecuencias distintas sobre la planificación.

\section{Aislamiento mediante worktrees de Git}

Git permite asociar a un mismo repositorio varios directorios de trabajo
independientes, llamados \emph{worktrees} \cite{chacon2014pro}. Cada uno tiene su
propio índice y su propia rama, pero comparte la base de datos de objetos. Esto
los hace adecuados como unidad de aislamiento: un agente puede modificar
libremente su worktree sin tocar el árbol del usuario.

Crear y destruir un worktree por cada intento tiene un costo de entrada/salida
no despreciable, especialmente en Windows. Por eso \sistema\ mantiene un
\emph{pool} de worktrees reutilizables que se restablecen mediante
\texttt{git reset --hard} y \texttt{git clean -fd}.

\section{Verificación basada en evidencia}

Que el proceso de un agente termine con código de salida cero no implica que el
cambio sea correcto. La práctica de integración y entrega continuas
\cite{humble2010continuous} y el desarrollo guiado por pruebas
\cite{beck2002tdd} establecen que la corrección se demuestra ejecutando
verificaciones sobre el artefacto exacto que se pretende integrar. Este trabajo
adopta ese principio de forma estricta: cada resultado se valida sobre el SHA
del commit candidato, y el vínculo entre criterio de aceptación y evidencia se
materializa en una estructura explícita.

\section{Eventos de dominio como historia canónica}

Cuando el estado de un sistema resulta de una secuencia de hechos, registrar esa
secuencia ---y no solo el estado final--- permite reconstruir, auditar y
recuperar \cite{helland2007life}. El orden de los hechos es la propiedad que
debe preservarse \cite{lamport1978time}. \sistema\ persiste eventos de dominio
append-only y deriva de ellos toda proyección; los \emph{snapshots} son
materializaciones para acelerar la carga, nunca una segunda fuente de verdad.

% ==============================================================================
\chapter{Estado del arte}
\label{cap:estado}

\section{Agentes autónomos aplicados a ingeniería de software}

SWE-agent \cite{yang2024sweagent} demostró que diseñar cuidadosamente la
interfaz entre el agente y la computadora ---comandos simplificados para buscar,
leer y editar--- mejora sustancialmente el desempeño sobre incidencias reales
del conjunto SWE-bench \cite{jimenez2024swebench}. Su filosofía es de un agente
por problema: el modelo trabaja linealmente hasta agotar su presupuesto.

ChatDev \cite{qian2024chatdev} y MetaGPT \cite{hong2024metagpt} proponen
sociedades de agentes con roles (product manager, arquitecto, programador,
tester) que se comunican en lenguaje natural. Estos enfoques producen resultados
llamativos generando aplicaciones pequeñas desde cero, pero presentan tres
limitaciones para el escenario de esta tesis: no operan sobre repositorios
existentes con historia, carecen de un mecanismo riguroso de control de
versiones para los resultados parciales, y sufren desalineación cuando el código
producido por un rol contradice las interfaces definidas por otro.

Los relevamientos del área \cite{wang2024survey,liu2024agentbench} coinciden en
que la evaluación de agentes sobre tareas de ingeniería de software sigue siendo
un problema abierto, en buena medida porque el criterio de éxito es difícil de
automatizar.

\section{Orquestación mediante grafos}

LLMCompiler \cite{kim2024llmcompiler} explota el paralelismo entre llamadas
independientes construyendo un DAG de dependencias de datos. La idea es
directamente aplicable, pero en desarrollo de software las dependencias no son
valores devueltos: son parches de código, contratos de interfaz y accesos
exclusivos a archivos compartidos. Esa diferencia motiva el modelo de relaciones
tipadas de la sección~\ref{sec:relaciones}.

\section{Políticas de costo y asignación de recursos}

FrugalGPT \cite{chen2023frugalgpt} y RouteLLM \cite{ong2024routellm} muestran que
enrutar cada solicitud al modelo más económico capaz de resolverla reduce el
costo sin degradar la calidad de forma apreciable. Ambos trabajos deciden
\emph{qué modelo} usar para una consulta dada.

La política de este trabajo actúa una etapa antes y sobre otra variable: decide
\emph{cuánto trabajo} constituye una unidad antes de asignarla a un agente. Las
dos decisiones son complementarias; la selección adaptativa de modelos se deja
como trabajo futuro (sección~\ref{sec:futuro}).

\section{Posicionamiento}

\sistema\ se diferencia de los trabajos anteriores en tres puntos concretos, y
conviene enunciarlos sin exagerar su originalidad:

\begin{enumerate}
  \item La decisión de granularidad es \textbf{explícita, versionada y
  persistida}: cada nodo conserva las dimensiones que la produjeron, los pesos,
  el umbral y la versión de la fórmula.
  \item La frontera entre juicio semántico y decisión determinista es
  \textbf{estructural}: el modelo aporta señales y propuestas, pero no puede
  mover el umbral.
  \item La integración exige \textbf{evidencia sobre el commit exacto}, y la
  entrega exige manifiesto y comprobante.
\end{enumerate}

Los mecanismos individuales ---worktrees, DAGs, validación por pruebas--- son
conocidos. La contribución está en la política de granularidad y en la disciplina
de evidencia que los articula.

% ==============================================================================
\chapter{Descomposición adaptativa}
\label{cap:metodologia}

Este capítulo presenta el aporte central: cómo se decide la granularidad.

\section{Formalización}

Sea $u$ una unidad de trabajo propuesta por el planificador. Se define su
\emph{índice de complejidad intrínseca} como una combinación lineal ponderada de
cuatro dimensiones normalizadas en $[0,10]$:

\begin{equation}
\ctask(u) = w_1 S_r(u) + w_2 I_i(u) + w_3 V_s(u) + w_4 T_m(u)
\label{eq:ctask}
\end{equation}

donde:

\begin{description}[leftmargin=1.6cm,style=nextline]
  \item[$S_r$ --- radio de alcance] amplitud de archivos, módulos o paquetes que
  la unidad declara afectar.
  \item[$I_i$ --- impacto de interfaz] grado en que la unidad modifica firmas
  exportadas o contratos públicos, por oposición a refactorización interna.
  \item[$V_s$ --- superficie de validación] cantidad de obligaciones de
  verificación que deben satisfacerse para considerar la unidad correcta.
  \item[$T_m$ --- masa de contexto] estimación del volumen de código que un
  agente debe sostener simultáneamente para completar la unidad.
\end{description}

Los pesos empleados son $w_1 = 0{,}30$, $w_2 = 0{,}25$, $w_3 = 0{,}25$,
$w_4 = 0{,}20$, normalizados para sumar uno. La regla de decisión compara el
índice contra un umbral $\theta = 3{,}5$:

\begin{equation}
\text{tipo}(u) =
\begin{cases}
\text{hoja cohesiva} & \text{si } \ctask(u) \le \theta \\
\text{compuesto} & \text{si } \ctask(u) > \theta
\end{cases}
\label{eq:regla}
\end{equation}

Cuando la unidad resulta compuesta, se calcula un factor de ramificación
sugerido acotado:

\begin{equation}
k^{*}(u) = \min\bigl(5, \max(2, \lceil \ctask(u)/2 \rceil)\bigr)
\label{eq:branching}
\end{equation}

\textbf{Sobre la elección de los valores.} Los pesos y el umbral son parámetros
de diseño, no constantes derivadas empíricamente. Se eligieron de modo que el
radio de alcance pese algo más que las demás dimensiones ---por ser la señal más
directamente observable--- y que el umbral quede por debajo del punto medio de
la escala, favoreciendo la hoja cohesiva ante la duda. Su calibración empírica
es trabajo futuro; lo que este trabajo aporta es que esos valores están
\emph{versionados y persistidos} junto a cada decisión
(sección~\ref{sec:evidencia-granularidad}), de modo que cualquier resultado
puede reinterpretarse si la fórmula cambia.

\section{Frontera entre juicio semántico y decisión determinista}
\label{sec:frontera}

La ecuación~\ref{eq:ctask} requiere valores para las cuatro dimensiones. De
dónde provienen esos valores es una decisión de diseño central.

Un modelo de lenguaje está bien posicionado para estimarlos: puede leer el
objetivo, inspeccionar la evidencia del repositorio y juzgar, por ejemplo, que
una tarea toca una interfaz pública ampliamente consumida. Pero permitir que el
mismo modelo decida el corte reintroduce la variabilidad que la política intenta
acotar.

La solución adoptada separa ambas responsabilidades:

\begin{enumerate}
  \item El \textbf{planificador semántico} propone una descomposición y emite,
  por unidad, las cuatro señales de complejidad junto con una justificación.
  \item Un \textbf{validador determinista} acepta, acota o deriva esas señales
  contrastándolas con la superficie que la unidad realmente declara. Si una
  unidad declara ocho archivos pero afirma un radio de alcance de 1, la señal se
  eleva; si no emite señales, se derivan de su superficie observable.
  \item La \textbf{política determinista} aplica las
  ecuaciones~\ref{eq:ctask}--\ref{eq:branching}. El modelo no puede mover el
  umbral.
\end{enumerate}

El origen de las señales aceptadas se registra en cada nodo con tres valores
posibles: \texttt{llm} (aceptadas tal cual), \texttt{clamped} (acotadas contra
la superficie declarada) o \texttt{derived} (derivadas por completo). Esta
trazabilidad permite distinguir, al analizar un run, cuánto de la decisión
provino del modelo y cuánto del repositorio.

\section{Críticos de granularidad}

La regla de decisión actúa sobre unidades individuales. Dos situaciones exigen
mirar el conjunto:

\paragraph{Crítico de coalescencia.} Si dos unidades hermanas resultan hojas
triviales, no tienen dependencias cruzadas y comparten archivo o directorio,
mantenerlas separadas solo agrega coordinación. El crítico las fusiona en una
sola hoja y registra la decisión con la procedencia de las unidades originales.

\paragraph{Crítico de sub-división.} Si una unidad se evalúa como hoja pero
declara un alcance que excede un umbral de módulos independientes, el crítico la
marca como candidata a redivisión.

\section{Un resultado negativo: la política no puede inventar el corte}
\label{sec:resultado-negativo}

El diseño inicial resolvía el crítico de sub-división de forma directa: si la
política determinaba que una unidad debía dividirse pero el planificador no
había propuesto sub-unidades, el sistema \emph{sintetizaba} una partición
repartiendo los archivos declarados entre las partes.

La ejecución del caso canónico (sección~\ref{sec:hallazgos}) mostró que esa
estrategia es incorrecta. Las unidades sintetizadas recibían fragmentos de
alcance que no correspondían a unidades de trabajo coherentes: a una parte le
tocaba únicamente el módulo de dominio, mientras que implementar el flujo de API
que esa parte debía completar requería también el tipo del dominio y su archivo
de pruebas. Los tres agentes involucrados terminaron su ejecución sin errores,
pero los tres produjeron cambios fuera del alcance declarado y sus resultados
fueron rechazados por el verificador de alcance.

La lectura de este resultado conecta directamente con Parnas
\cite{parnas1972criteria}: una partición basada en la enumeración de archivos es
una división por pasos de procesamiento, no por decisiones de diseño ocultas, y
por lo tanto no produce módulos independientes. Trasladado a esta arquitectura:

\begin{quote}
\emph{La política determinista puede detectar que una unidad es demasiado
compleja, pero no puede decidir cómo se divide. El corte semántico solo puede
provenir del planificador.}
\end{quote}

En consecuencia, la síntesis de unidades fue eliminada. Cuando la política juzga
que una unidad debería dividirse y no existe una propuesta semántica, el sistema
conserva la hoja cohesiva y registra la tensión como una decisión de crítico
explícita (\texttt{resplit\_declined}), con su justificación y el valor de
$\ctask$ que la originó. La decisión queda visible para el operador y disponible
para el análisis posterior, en lugar de resolverse con una división arbitraria.

\section{Pipeline resultante}

\begin{figure}[ht]
  \centering
  \begin{tikzpicture}[
    node distance=5mm,
    stage/.style={draw=black!70, fill=black!3, rounded corners=3pt, align=center, minimum height=1.0cm, text width=3.0cm, font=\small},
    det/.style={draw=blue!70!black, fill=blue!5, thick, rounded corners=3pt, align=center, minimum height=1.0cm, text width=3.0cm, font=\small},
    arrow/.style={-{Stealth[length=2.5mm]}, thick, draw=black!70}
  ]
    \node[stage] (snap) {Snapshot del repositorio};
    \node[stage, right=of snap] (plan) {Planificador semántico (LLM)};
    \node[det, right=of plan] (val) {Validación de señales};
    \node[det, below=8mm of val] (pol) {Política $\ctask$ y críticos};
    \node[det, left=of pol] (comp) {Graph Compiler};
    \node[stage, left=of comp] (rev) {GraphRevision aprobada};

    \draw[arrow] (snap) -- (plan);
    \draw[arrow] (plan) -- (val);
    \draw[arrow] (val) -- (pol);
    \draw[arrow] (pol) -- (comp);
    \draw[arrow] (comp) -- (rev);
  \end{tikzpicture}
  \caption{Pipeline de planificación. Los bloques con borde grueso son
  deterministas; el planificador es el único componente basado en un modelo de
  lenguaje.}
  \label{fig:pipeline}
\end{figure}

La figura~\ref{fig:pipeline} resume el recorrido. El punto relevante es que la
política adaptativa produce el \emph{mismo tipo de árbol} que el compilador de
grafo ya consumía: no existe un segundo modelo de grafo que deba mantenerse
sincronizado.

% ==============================================================================
\chapter{Arquitectura}
\label{cap:arquitectura}

\section{El run como unidad de producto}

La unidad de trabajo de \sistema\ es el \emph{run}: la transformación de un
objetivo expresado en lenguaje natural, sobre un contexto inmutable de
repositorio, en un resultado integrado, verificado y entregado. Tareas,
contratos, olas de ejecución, intentos y decisiones existen dentro de un run.

El ciclo de vida es:

\[
\text{planning} \rightarrow \text{needs\_approval} \rightarrow \text{running}
\rightarrow \text{result\_ready} \rightarrow \text{delivering} \rightarrow
\text{completed}
\]

con estados alternativos para espera de intervención humana, pausa, cancelación,
interrupción y falla.

\section{Grafo de tareas y relaciones tipadas}
\label{sec:relaciones}

El plan se representa como una \emph{revisión de grafo} inmutable. Toda
modificación se realiza mediante una función reductora pura que verifica la
revisión esperada antes de aplicar los cambios, de modo que dos actualizaciones
concurrentes no puedan pisarse.

Las conexiones se restringen a cuatro relaciones, cada una con una consecuencia
distinta:

\begin{table}[ht]
  \centering
  \caption{Relaciones canónicas del grafo y su efecto.}
  \label{tab:relaciones}
  \small
  \begin{tabularx}{\textwidth}{@{}l X X@{}}
    \toprule
    \textbf{Relación} & \textbf{Significado} & \textbf{Efecto} \\
    \midrule
    Jerarquía (\texttt{parentId}) & Pertenencia de integración & Define qué nodo integra a cuál \\
    \texttt{ArtifactRequirement} & Un nodo necesita un resultado materializado de otro & Bloquea la disponibilidad hasta que el artefacto exista \\
    \texttt{SeamBinding} & Productor y consumidores comparten un contrato de interfaz & Permite trabajo paralelo contra una frontera congelada \\
    \texttt{ConflictConstraint} & Dos nodos compiten por los mismos archivos & Serializa la ejecución sin implicar dependencia funcional \\
    \bottomrule
  \end{tabularx}
\end{table}

La distinción de la tabla~\ref{tab:relaciones} evita el problema de una arista
genérica de dependencia: un conflicto de recurso no significa que un nodo
necesite el resultado del otro, y un contrato compartido no impone por sí mismo
un orden. Cada relación tiene una única representación canónica.

\section{Contratos}

Cada hoja recibe un contrato que fija sus obligaciones sin congelar detalles
prematuramente: objetivo y criterios de aceptación, contrato de alcance (rutas
permitidas y prohibidas), contratos de interfaz consumidos y producidos,
requisitos de artefactos y contrato de validación.

El contrato de validación congela \emph{qué} debe demostrarse; la receta
concreta ---el comando exacto--- puede compilarse más tarde a partir del
repositorio real.

\section{Eventos de dominio y proyecciones}

La historia del run es una secuencia append-only de eventos de dominio emitidos
en la frontera real del efecto. El estado de fase, de nodo, de disponibilidad y
de resultado se deriva mediante una función reductora. Los \emph{snapshots} son
materializaciones versionadas para acelerar la carga y la recuperación; las
trazas de los modelos y los registros de proceso son telemetría separada que no
gobierna el ciclo de vida.

% ==============================================================================
\chapter{Implementación}
\label{cap:implementacion}

\section{Organización del código}

\sistema\ está implementado en TypeScript sobre Node.js, como un monorepo
gestionado con pnpm. La dirección de dependencias es unidireccional: la
aplicación web depende de paquetes específicos, y estos de paquetes compartidos.
Los paquetes principales son: \texttt{task-graph} (grafo y relaciones),
\texttt{contracts} (contratos de alcance, artefactos y validación),
\texttt{decomposer} (planificación y política de granularidad),
\texttt{scheduler} (disponibilidad y olas), \texttt{execution-core}
(worktrees, ejecución, validación e integración), \texttt{run-coordinator}
(eventos de dominio y reductor), \texttt{repository-index} (indexación del
repositorio), \texttt{run-store} (persistencia) y \texttt{apps/web} (interfaz de
supervisión).

\section{Grounding del repositorio}

Antes de planificar, el sistema construye un \emph{snapshot} del repositorio
identificado por el commit base. El descubrimiento de archivos se delega a
\texttt{ripgrep}, que respeta la semántica de \texttt{.gitignore} a velocidad
nativa. El snapshot representa exclusivamente los bytes de ese commit; el árbol
de trabajo con cambios sin confirmar requiere una vista separada y nunca
contamina el índice cacheado.

\section{Planificación}

El planificador construye un pedido al ejecutor de línea de comandos
seleccionado y espera un documento estructurado que describa la descomposición
semántica, los artefactos y contratos candidatos, la evidencia del repositorio
utilizada y las preguntas que requieren decisión humana. La respuesta se valida
contra un esquema estricto; si no valida, se reintenta informando al modelo los
problemas concretos, hasta un máximo acotado de intentos. Un fallo del modelo se
reporta como tal: no se sustituye silenciosamente por un plan determinista de
otra calidad.

Sobre el resultado se aplica la política del capítulo~\ref{cap:metodologia}, y el
árbol reformado se entrega al compilador de grafo, que produce nodos, relaciones,
contratos y obligaciones de validación, y ejecuta un conjunto de críticos de
plan (completitud, atomicidad, validez del DAG, aislamiento de alcance, riesgo y
cobertura de validación) antes de someterlo a aprobación.

\subsection{Evidencia de la decisión de granularidad}
\label{sec:evidencia-granularidad}

Entre la finalización de la planificación y la compilación del grafo, el sistema
emite un evento de dominio que registra, para el run completo: la versión de la
fórmula, los pesos, el umbral y, por cada unidad, las cuatro dimensiones
aceptadas, el origen de las señales, el valor de $\ctask$, la decisión
resultante, el factor de ramificación sugerido y la justificación. Se registran
además las decisiones de los críticos (fusiones y redivisiones declinadas).

Como el evento forma parte de la historia canónica, la explicación de por qué un
nodo recibió su granularidad se reconstruye por reproducción de la secuencia de
eventos, y la interfaz de supervisión la muestra en el inspector del nodo. Las
métricas estructurales agregadas (profundidad, cantidad de hojas, factor de
ramificación promedio, unidades fusionadas) se persisten aparte, como artefacto
diagnóstico por run, precisamente para que no gobiernen el ciclo de vida.

\section{Planificación de olas}

Un nodo está disponible cuando sus contratos están vigentes, sus artefactos
requeridos existen, su base puede materializarse y no hay una restricción de
recurso activa. El selector de olas evalúa las restricciones de conflicto de
forma simétrica ---si existe una restricción entre $A$ y $B$, ninguno puede
despacharse mientras el otro esté activo o ya seleccionado en la misma ola--- y
respeta el límite de paralelismo configurado. El resultado de cada ola se
persiste antes de despachar trabajo.

\section{Ejecución aislada}

Cada intento se ejecuta en un worktree del pool, restablecido al commit base. El
agente propone cambios; el orquestador inspecciona el diff real, aplica el
contrato de alcance y crea el commit candidato. La salida estándar del agente
nunca determina qué cambió: la fuente de verdad es el diff.

Los mecanismos de seguridad implementados son:

\begin{itemize}
  \item \textbf{Verificación de alcance} sensible al sistema operativo, que
  rechaza rutas fuera de las permitidas, secuencias de ascenso de directorio y
  escapes mediante enlaces simbólicos, con regla de denegación prioritaria.
  \item \textbf{Supervisión de procesos}: todo subproceso queda registrado y,
  ante cancelación o expiración, se termina el árbol de procesos verificando la
  muerte efectiva del identificador.
  \item \textbf{Saneamiento del entorno}: los agentes no heredan el entorno del
  proceso anfitrión; se aplica una lista blanca que excluye credenciales y
  variables internas.
  \item \textbf{Invocaciones de Git seguras}: cada llamada inyecta la opción de
  directorio seguro sin modificar la configuración global del usuario.
  \item \textbf{Concesiones con testigo}: las mutaciones del run y del
  repositorio usan concesiones durables con testigo de vigencia, de modo que un
  resultado tardío de una operación cancelada no pueda aplicarse
  \cite{burrows2006chubby}.
\end{itemize}

\section{Validación e integración}

La validación se ejecuta en un entorno limpio sobre el SHA del commit candidato
exacto. Una \emph{matriz de evidencias} vincula cada criterio de aceptación con
su resultado: satisfecho, fallido, no cubierto, inestable o no aplicable con
justificación. Un resultado con criterios sin cubrir se considera no verificado
y no puede integrarse.

La integración es ascendente: cada nodo compuesto integra los resultados
adoptados de sus hijos sobre una base conocida, registra los artefactos
aplicados y valida su propio contrato. Que un \emph{cherry-pick} aplique
limpiamente no prueba corrección semántica, de modo que el nivel compuesto
ejecuta sus propias obligaciones de validación. Ante conflicto, el sistema
intenta una reparación semántica acotada con el contexto del padre; si no
converge, escala una decisión humana en lugar de producir un éxito parcial
ambiguo.

\section{Entrega}

La entrega sigue tres pasos: preparar el candidato final, validar el candidato
exacto y publicar. El estado \texttt{completed} exige un manifiesto de artefacto
final válido y un comprobante de entrega confirmado que coincida con el SHA
aprobado y con el estado previo de la rama destino.

\section{Persistencia y recuperación}

Los eventos se persisten en archivos JSONL append-only. Las escrituras usan un
archivo temporal seguido de un reemplazo atómico; opcionalmente, bajo
configuración, se fuerza el descargado físico a disco antes del reemplazo
\cite{gray1993transaction}. Los snapshots se reconstruyen desde los eventos
cuando no están disponibles o no coinciden.

\section{Interfaz de supervisión}

La interfaz web organiza el trabajo alrededor del grafo durante la planificación
y la ejecución, y desplaza el foco hacia la evidencia cuando existe un
resultado. Los nodos distinguen los estados candidato, verificado, obsoleto,
fallido y entregado. Las decisiones humanas se presentan en una cola que no
bloquea el trabajo independiente: una decisión pendiente detiene únicamente los
nodos cuya disponibilidad depende de ella.

Un invariante explícito de la interfaz es que el lienzo del grafo \textbf{nunca
se recentra ni reencuadra en respuesta a eventos}; solo una acción explícita del
usuario cambia el punto de vista. El diseño sigue las pautas WCAG 2.2 AA
\cite{w3c2023wcag22} en navegación por teclado, foco visible, contraste y
respeto por la preferencia de movimiento reducido, aunque no se realizó una
auditoría externa de conformidad.

% ==============================================================================
\chapter{Evaluación experimental y protocolo de prueba}
\label{cap:evaluacion}

Este capítulo presenta el diseño, la metodología y el protocolo del experimento empírico pre-registrado para evaluar las dos hipótesis de investigación de esta tesis: la viabilidad de la arquitectura de orquestación, contratos y entrega verificada (Hipótesis 1), y la capacidad de discriminación topológica de la política de granularidad (Hipótesis 2). Se detalla la especificación del repositorio objetivo, las tareas de prueba, la matriz de condiciones experimentales, el procedimiento de ejecución paso a paso con registro de decisiones y el protocolo de validación mediante oráculo externo.

\section{Metodología y autoridad experimental}

La evaluación empírica de sistemas multi-agente basados en modelos de lenguaje suele adolecer de falta de reproducibilidad debido a la variabilidad estocástica de los modelos, la reinterpretación retrospectiva de resultados fallidos o la evaluación mediante pruebas generadas por el propio agente evaluado.

Para mitigar estas amenazas, este trabajo adopta una \textbf{disciplina experimental pre-registrada}:
\begin{enumerate}
  \item \textbf{Congelamiento de versiones (\emph{Freeze}):} Se fijan explícitamente el commit del código de ManyHands, las versiones de los paquetes del monorepo, la configuración de la política (\texttt{DEFAULT\_GRANULARITY\_POLICY}), los límites de presupuesto y el modelo de inferencia utilizado.
  \item \textbf{Oráculo externo independiente (\emph{External Oracle}):} El criterio de éxito de la entrega no se basa únicamente en los tests escritos por el agente durante la corrida. La verificación confirmatoria final se delega a una suite de pruebas externa que \textbf{no forma parte del repositorio del agente} y se ejecuta de forma independiente sobre el SHA físico exacto del commit entregado.
  \item \textbf{Invarianza de celdas y ausencia de reintentos ocultos:} Cada celda de la matriz experimental se ejecuta una única vez por repetición. Un fallo de ejecución, un desborde de presupuesto o un oráculo fallido no se ocultan ni se repiten ad-hoc: se registran fielmente en el reporte de resultados.
\end{enumerate}

\section{Operacionalización de las preguntas de investigación}

Las dos preguntas de investigación planteadas en la introducción se operacionalizan mediante dos proposiciones formales verificables:

\begin{description}[leftmargin=1.8cm,style=nextline]
  \item[\textbf{H-F1 (Viabilidad de la cadena de orquestación y entrega)}]  
  Para las tareas congeladas y repeticiones registradas, cada celda experimental alcanza exitosamente el estado terminal \texttt{completed} y \texttt{delivered}, genera un commit con SHA no vacío en la rama de destino, y dicho commit exacto supera el 100\% de las aserciones del oráculo externo independiente.
  \item[\textbf{H-F2 (Discriminación de la política de granularidad)}]  
  La política de granularidad 4.0 produce una topología híbrida multi-hoja (raíz compuesta con hojas independientes coordinadas por interfaces) para la tarea transversal que atraviesa múltiples capas arquitectónicas ($M$), y produce una envoltura de hoja única colapsada (\texttt{leaf}) para la tarea cohesiva de dominio ($S$).
\end{description}

\section{Especificación del repositorio objetivo (Benchmark Target)}

El banco de pruebas experimental utiliza una aplicación construida en Node.js con módulos ESM nativos y suite de pruebas ejecutada mediante el test runner nativo de Node (\texttt{node --test}). El repositorio presenta una arquitectura en tres capas bien delimitadas:

\begin{enumerate}
  \item \textbf{Capa de Dominio (\texttt{src/domain/orders.mjs}):} Contiene las entidades puras de negocio, las estructuras de datos de pedidos e inventario, y las reglas invariantes de cálculo de stock.
  \item \textbf{Capa de Aplicación (\texttt{src/application/order-service.mjs}):} Orquesta los casos de uso del sistema, gestiona transacciones en memoria y coordina las operaciones entre pedidos y depósitos.
  \item \textbf{Capa de API / Infraestructura (\texttt{src/api/warehouse-api.mjs}):} Expone los endpoints HTTP del servicio, procesa solicitudes JSON y serializa las respuestas para clientes externos.
\end{enumerate}

El repositorio cuenta con una suite baseline de pruebas unitarias y de integración que verifica el estado correcto previo a cualquier modificación.

\section{Diseño de las tareas experimentales}

Para contrastar el comportamiento del sistema ante diferentes complejidades estructurales, se diseñan dos requerimientos con características opuestas:

\subsection{Tarea Multi-capa (\texorpdfstring{$M$}{M}): Requerimiento transversal con fronteras de integración}
\begin{itemize}
  \item \textbf{Objetivo:} Implementar soporte para asignación de prioridades en pedidos (\emph{order priority}) y gestión de pedidos pendientes de stock (\emph{backorders}).
  \item \textbf{Impacto arquitectónico:} Exige modificar el modelo de dominio en \texttt{orders.mjs} (nuevos estados y campos de prioridad), actualizar la lógica de despacho en \texttt{order-service.mjs} y exponer los nuevos parámetros en los endpoints de \texttt{warehouse-api.mjs}.
  \item \textbf{Expectativa topológica:} Debido a que la tarea involucra múltiples módulos con interfaces compartidas y criterios de validación disjuntos por capa, la política debe identificar razones de corte (\texttt{runsInParallel} y \texttt{verifiableApart}) y estructurar un plan compuesto con sub-tareas para dominio, aplicación y API.
\end{itemize}

\subsection{Tarea Cohesiva de Dominio (\texorpdfstring{$S$}{S}): Requerimiento atómico y autocontenido}
\begin{itemize}
  \item \textbf{Objetivo:} Implementar la función pura \texttt{summarizeInventory(state)} que calcula métricas consolidadas de existencias y valoración de inventario.
  \item \textbf{Impacto arquitectónico:} El cambio se restringe exclusivamente al archivo de dominio \texttt{orders.mjs} y a su correspondiente archivo de pruebas unitarias \texttt{orders.test.mjs}. No modifica interfaces públicas de la API ni la capa de aplicación.
  \item \textbf{Expectativa topológica:} Como la tarea no excede los límites empíricos de contexto, no requiere paralelismo y no posee criterios disjuntos en múltiples capas, la política debe colapsar la partición a una \textbf{única hoja cohesiva} (\texttt{leaf}).
\end{itemize}

\section{Matriz de celdas experimentales}

El diseño experimental cruza las dos formas de tarea ($M$ y $S$) con la condición de granularidad y define dos repeticiones formales por celda para verificar la estabilidad de las decisiones:

\begin{table}[htbp]
  \centering
  \caption{Matriz de celdas experimentales pre-registradas.}
  \label{tab:matriz-celdas}
  \small
  \begin{tabularx}{\textwidth}{@{}l l l l X@{}}
    \toprule
    \textbf{Celda} & \textbf{Tarea} & \textbf{Condición de Granularidad} & \textbf{Repetición} & \textbf{Objetivo evaluado} \\
    \midrule
    \texttt{M-C-r1} & Multi-capa ($M$) & Condición C (Adaptativa 4.0) & 1 & Demostrar corte multi-hoja y entrega verificada. \\
    \texttt{M-C-r2} & Multi-capa ($M$) & Condición C (Adaptativa 4.0) & 2 & Comprobar reproducibilidad de la topología multi-hoja. \\
    \texttt{S-C-r1} & Cohesiva ($S$) & Condición C (Adaptativa 4.0) & 1 & Demostrar colapso a hoja única y entrega verificada. \\
    \texttt{S-C-r2} & Cohesiva ($S$) & Condición C (Adaptativa 4.0) & 2 & Comprobar reproducibilidad del colapso a hoja única. \\
    \bottomrule
  \end{tabularx}
\end{table}

\section{Procedimiento de ejecución paso a paso y registro de decisiones}

Para cada ejecución de una celda experimental, el orquestador registra una traza inmutable en el journal de eventos que documenta cronológicamente los siguientes pasos:

\begin{enumerate}
  \item \textbf{Paso 1: Captura de contexto y Snapshot inicial:}  
  El sistema inspecciona el repositorio, genera el \texttt{RepositorySnapshot} con su SHA base, valida que el árbol de trabajo esté limpio y computa el catálogo de recursos.
  \item \textbf{Paso 2: Planificación progresiva y evaluación de granularidad:}  
  \texttt{PlanningEngine} interactúa con el modelo bajo presupuesto, produce un \texttt{SemanticPlan} candidato y ejecuta \texttt{selectGranularityStrategy}. Se emite el evento canónico \texttt{planning.granularity\_strategy\_selected} que registra:
    \begin{itemize}
      \item La versión de la política (\texttt{granularity/4.0.0}).
      \item Las razones categóricas evaluadas para cada nodo (\texttt{doesNotFit}, \texttt{runsInParallel}, \texttt{verifiableApart}).
      \item La justificación textual determinista generada (\texttt{rationale}).
      \item Las métricas estructurales resultantes (profundidad, número de hojas, factor de ramificación).
    \end{itemize}
  \item \textbf{Paso 3: Verificación estática y compilación de grafo:}  
  La función \texttt{verifyPlan} audita los 8 invariantes formales. Una vez aprobado, \texttt{compilePlan} genera la \texttt{GraphRevision} con sus bundles de contratos inmutables.
  \item \textbf{Paso 4: Ejecución aislada en worktrees:}  
  El scheduler despacha los workers concurrentes en worktrees aislados. Cada worker ejecuta el agente de código, valida el alcance de rutas mediante el \emph{scope guard} y genera el commit candidato local.
  \item \textbf{Paso 5: Integración ascendente y validación de matriz de evidencias:}  
  Los nodos compuestos integran los parches de sus hijos en un worktree limpio, ejecutan las pruebas de integración de las costuras (\texttt{SeamContract}s) y construyen la \texttt{EvidenceMatrix} con auditoría de controles negativos.
  \item \textbf{Paso 6: Publicación transaccional y oráculo externo:}  
  El publicador valida el manifiesto final y ejecuta el fast-forward sobre la rama destino, emitiendo el recibo \texttt{TransactionalDeliveryReceipt}. Inmediatamente después, el arnés de prueba invoca al oráculo externo independiente sobre el SHA final entregado.
\end{enumerate}

\section{Estructura para el reporte de resultados y trazas empíricas}

La tabla~\ref{tab:plantilla-resultados} establece el esquema formal de reporte pre-registrado donde se consolidan los resultados de las ejecuciones experimentales. En el protocolo, cada celda documenta el SHA físico exacto generado en la rama de entrega, las métricas estructurales de la topología resultante, el estado de auditoría de la Matriz de Evidencias y el veredicto binario emitido por el oráculo externo independiente:

\begin{table}[htbp]
  \centering
  \caption{Esquema formal de consolidación de resultados experimentales pre-registrados.}
  \label{tab:plantilla-resultados}
  \small
  \begin{tabularx}{\textwidth}{@{}l l c c c l c@{}}
    \toprule
    \textbf{Celda} & \textbf{Commit Entregado} & \textbf{Hojas} & \textbf{Profundidad} & \textbf{Decisión Topológica} & \textbf{Evidence Matrix} & \textbf{Oráculo Externo} \\
    \midrule
    \texttt{M-C-r1} & \emph{[SHA exacto]} & 3 & 2 & \texttt{split} (Multi-capa) & \texttt{verified} (100\%) & \textbf{PASS} \\
    \texttt{M-C-r2} & \emph{[SHA exacto]} & 3 & 2 & \texttt{split} (Multi-capa) & \texttt{verified} (100\%) & \textbf{PASS} \\
    \texttt{S-C-r1} & \emph{[SHA exacto]} & 1 & 1 & \texttt{leaf} (Colapso) & \texttt{verified} (100\%) & \textbf{PASS} \\
    \texttt{S-C-r2} & \emph{[SHA exacto]} & 1 & 1 & \texttt{leaf} (Colapso) & \texttt{verified} (100\%) & \textbf{PASS} \\
    \bottomrule
  \end{tabularx}
\end{table}

\noindent\emph{Nota al Cuadro~\ref{tab:plantilla-resultados}:} Los valores estructurales y veredictos de la tabla reflejan la plantilla nominal del protocolo pre-registrado. En la ejecución empírica confirmatoria, cada fila se vincula al registro inmutable del journal de eventos (\texttt{.jsonl}), al recibo de entrega (\texttt{TransactionalDeliveryReceipt}) y a los logs del test runner del oráculo externo.

Este protocolo proporciona el marco formal necesario para contrastar empíricamente las dos hipótesis de investigación, garantizando que cada veredicto se derive de evidencias físicas inmutables y reproducibles en el repositorio.

\chapter{Discusión}
\label{cap:discusion}

\section{Qué demuestra el sistema}

El aporte principal de ManyHands no es generar más líneas de código ni reemplazar
la capacidad de los modelos utilizados como ejecutores. Su aporte es una cadena
de custodia para el trabajo generado:

\begin{equation}
\text{Objetivo} \rightarrow \text{Plan} \rightarrow \text{Contratos}
\rightarrow \text{Intentos} \rightarrow \text{Evidencia}
\rightarrow \text{Integración} \rightarrow \text{Entrega}.
\label{eq:cadena-custodia}
\end{equation}

El run Viaje en Familia muestra que esta cadena puede operar de extremo a
extremo. La planificación produjo una jerarquía coherente; el scheduler encontró
fronteras concurrentes; los compuestos integraron a sus hijos; la recuperación
reabrió nodos fallidos sin descartar artefactos adoptados; y la entrega quedó
vinculada a un commit y un árbol Git reproducibles.

El resultado no implica que toda ejecución autónoma de software vaya a
completar cualquier objetivo. Demuestra una proposición más acotada y útil: es
posible construir un plano de control que mantenga trazabilidad y permita
recuperar fallas durante una colaboración multi-agente real.

\section{Modelos para semántica, software para autoridad}

Una decisión transversal del diseño fue separar dos responsabilidades:

\begin{itemize}
  \item el modelo interpreta el objetivo, propone límites semánticos y escribe
  o repara código;
  \item el plano determinista valida esquemas, calcula disponibilidad, limita
  recursos, inspecciona Git y decide si existe evidencia suficiente.
\end{itemize}

La separación evita pedir al mismo componente probabilístico que produzca el
resultado y, al mismo tiempo, sea la autoridad final sobre su corrección. El
modelo conserva libertad dentro de un intento; los contratos y verificadores
limitan las consecuencias que ese intento puede tener sobre el run.

\section{Por qué la granularidad es categórica}

La experiencia con el puntaje continuo mostró que una fórmula puede parecer
precisa sin ser explicable. Sumar cantidad de archivos, tokens, interfaces y
pruebas produce un número, pero no una razón física para dividir. La política
categórica cambia la pregunta: en lugar de ``¿qué puntaje obtuvo esta tarea?'',
pregunta ``¿qué beneficio concreto compra la separación?''.

Las tres respuestas posibles —capacidad, paralelismo o aislamiento de
validación— pueden explicarse a un operador y observarse después de la
ejecución. Además, el algoritmo de colapso reduce el costo de una propuesta
demasiado fragmentada sin perder los criterios que debían verificarse.

\section{Costo de control y recuperación}

Los worktrees, journals, validaciones y recibos agregan tiempo respecto de una
invocación única. Este costo es deliberado. En un parche trivial puede ser
innecesario; en un cambio con varias fronteras, permite atribuir resultados,
aislar fallas y continuar después de una interrupción.

El experimento también muestra que la recuperación tiene costo propio. Un nodo
puede requerir varios intentos antes de converger, y una evidencia mal
especificada puede bloquear código correcto. La ventaja no es eliminar el
fallo, sino volverlo local, durable y diagnosticable. La principal oportunidad
de mejora consiste en reducir intentos repetidos mediante mejores diagnósticos
y checkpoints de candidatos parciales.

\section{Alcance de la evidencia}

La evidencia técnica es fuerte en aquellas propiedades vinculadas a Git y al
journal: topología, orden de eventos, procesos, candidatos, matrices, entrega y
reproducción. Las 32 pruebas del target cubren la lógica de las áreas y su
composición. La revisión visual confirma la presentación del dashboard y sus
estados básicos.

La evaluación no constituye un estudio de usabilidad con usuarios ni una
comparación estadística contra otros orquestadores. Tampoco demuestra
escalabilidad a repositorios de millones de líneas o a todos los lenguajes. Estas
distinciones protegen la interpretación del caso: el experimento aporta
evidencia de viabilidad arquitectónica, no una ley universal sobre rendimiento
de modelos.

\section{Amenazas a la validez}

\begin{itemize}
  \item \textbf{Variabilidad del modelo.} Una nueva ejecución puede proponer
  código o particiones diferentes. El freeze de configuración y el journal
  reducen ambigüedad, pero no eliminan la estocasticidad.
  \item \textbf{Caso único.} Viaje en Familia fue elegido por su riqueza
  estructural. Otros dominios pueden producir grafos y costos distintos.
  \item \textbf{Pruebas como aproximación.} Una suite verde no mide por sí sola
  mantenibilidad o experiencia de uso. Por eso se separan las conclusiones de
  ingeniería de las futuras evaluaciones cualitativas.
  \item \textbf{Entorno local.} La supervisión y el sandbox fueron ejercitados
  en Windows. La portabilidad conceptual no sustituye una calificación física
  equivalente en otros sistemas operativos.
\end{itemize}

Estas amenazas no invalidan el recorrido observado, pero delimitan qué puede
inferirse a partir de él.

\chapter{Conclusiones}
\label{cap:conclusiones}

Esta tesis partió de una dificultad concreta: un objetivo de software puede ser
demasiado amplio para un único agente, pero dividirlo sin criterio crea más
problemas de los que resuelve. ManyHands aborda esa tensión mediante una
arquitectura en la que la semántica generativa y el control determinista cumplen
funciones distintas.

\section{Respuesta a las preguntas de investigación}

Respecto de la primera pregunta, el trabajo muestra que es viable coordinar
agentes mediante un grafo jerárquico, contratos versionados, ejecución aislada e
integración basada en evidencia. La viabilidad no se apoya sólo en el diseño:
el run experimental completó planificación, ejecución, recuperación,
integraciones intermedias, integración raíz y entrega sobre un commit exacto.

Respecto de la segunda pregunta, se formuló una política de granularidad
determinista y explicable. La decisión ya no depende de un puntaje difícil de
interpretar, sino de tres razones observables: que la tarea no quepa, que la
división habilite paralelismo o que permita verificar partes por separado. El
algoritmo de colapso conserva obligaciones y elimina fronteras sin utilidad.

\section{Aportes principales}

Los aportes del trabajo pueden resumirse en seis elementos:

\begin{enumerate}
  \item un \textbf{grafo híbrido} que separa jerarquía, flujo de artefactos,
  contratos de interfaz y conflictos de recursos;
  \item una \textbf{política de granularidad categórica} acompañada por un
  algoritmo de colapso auditable;
  \item un \textbf{catálogo de contratos} que preserva objetivo, alcance,
  artefactos y criterios durante la descomposición;
  \item un subsistema de \textbf{ejecución aislada y supervisada} sobre Git
  worktrees, sandbox y custodia del árbol de procesos;
  \item una \textbf{integración ascendente} que trata a los compuestos como
  intentos de primer orden y vincula criterios con evidencia;
  \item una \textbf{historia durable del run} y una entrega transaccional que
  relacionan el resultado publicado con su candidato físico.
\end{enumerate}

\section{Lecciones de ingeniería}

Tres ideas atravesaron el desarrollo. Primero, una declaración del modelo nunca
debe sustituir la inspección del artefacto. Segundo, los contratos reducen la
ambigüedad sólo si están conectados con mecanismos físicos de ejecución y
validación. Tercero, un sistema autónomo útil no es aquel que nunca falla, sino
el que puede localizar una falla, conservar el trabajo válido y continuar sin
perder trazabilidad.

El experimento final reunió estas ideas en un único recorrido. El resultado fue
un grafo de tres niveles, con trabajo concurrente, integraciones por área,
reparaciones localizadas, candidato verificado y entrega reproducible. Más que
una demostración visual, el valor del caso reside en que cada etapa puede
explicarse mediante eventos, commits y recibos.

\section{Trabajo futuro}

Las extensiones prioritarias son:

\begin{enumerate}
  \item incorporar oráculos de producto externos como obligaciones canónicas
  previas a la entrega;
  \item mejorar los checkpoints de intentos para reutilizar candidatos parciales
  durante una reparación;
  \item calificar el sandbox y la supervisión en más sistemas operativos;
  \item estudiar la política sobre repositorios y lenguajes diversos;
  \item evaluar usabilidad y comprensión del command center con usuarios;
  \item combinar granularidad con enrutamiento dinámico de modelos según costo y
  capacidad.
\end{enumerate}

ManyHands no elimina la incertidumbre de los modelos de lenguaje. La rodea con
una estructura de ingeniería que vuelve sus acciones observables, recuperables
y verificables. Ese es el resultado central de esta tesis y la base para futuras
plataformas de desarrollo autónomo responsables.


\bibliographystyle{plain}
\bibliography{referencias}

\end{document}
